\section{The register map}
\label{sec:regmap}

Thirteen registers, each of them 32 bits. \textbf{The index} is what goes out on the wire; the
offset column is always $\text{index} \times 4$ and is a legacy of an older, memory-mapped variant
of the same design.

\begin{booktable}{@{}rlllL@{}}
\thead{Idx} & \thead{Register} & \thead{Offset} & \thead{Acc.} & \thead{Description} \\ \midrule
0 & \code{STATUS} & \code{0x00} & R & Bit 0: TX ready. Bit 1: RX valid. Bit 2: Error. \\
1 & \code{TX_ID} & \code{0x04} & R/W & 11-bit transmit ID (bits 10--0). \\
2 & \code{TX_DLC} & \code{0x08} & R/W & Data Length Code (bits 3--0, value 0--8). \\
3 & \code{TX_DATA_LO} & \code{0x0C} & R/W & Transmit data bytes 0--3 (byte 0 = MSB). \\
4 & \code{TX_DATA_HI} & \code{0x10} & R/W & Transmit data bytes 4--7. \\
5 & \code{TX_SEND} & \code{0x14} & W & Write \code{0x1} to trigger a transmission. \\
6 & \code{RX_ID} & \code{0x18} & R & Received identifier. \\
7 & \code{RX_DLC} & \code{0x1C} & R & Received DLC. \\
8 & \code{RX_DATA_LO} & \code{0x20} & R & Received data bytes 0--3. \\
9 & \code{RX_DATA_HI} & \code{0x24} & R & Received data bytes 4--7. \\
10 & \code{RX_ACK} & \code{0x28} & W & Write \code{0x1} to acknowledge. \\
11 & \code{ERROR_FLAGS} & \code{0x2C} & R/W & The error register; write \code{0x0} to clear. \\
12 & \code{TX_ABORT} & \code{0x30} & W & Write \code{0x1} to force TX ready back. \\
\end{booktable}

Indices 13--15 are reserved: a read gives \code{0x00000000}, a write is ignored.

\subsection{STATUS: the three bits everything revolves around}

\begin{booktable}{@{}rlLL@{}}
\thead{Bit} & \thead{Meaning} & \thead{Set by} & \thead{Cleared by} \\ \midrule
0 & TX ready & reset; a completed transmission; a rising edge on the error level while the bit is
  low; a write to \code{TX_ABORT} & an accepted write to \code{TX_SEND} \\
1 & RX valid & a received frame, which latches \code{RX_*} at the same time & a write of
  \code{0x1} to \code{RX_ACK} \\
2 & Error & a \textbf{rising edge} on the controller's error level & a write of \code{0x0} to
  \code{ERROR_FLAGS} \\
\end{booktable}

Bits 31--3 always read as \code{0}.

\textbf{Read that table once more.} Every bit has an explicit clear, and that clearing is the
driver's responsibility. Forget \code{RX_ACK} and bit 1 stays set forever and you read the same
frame in an endless loop. Forget that \code{TX_SEND} clears bit 0 and you will wonder why the
second transmission is never accepted.

\subsection{Rules that are easy to miss}

\begin{itemize}
  \item \code{TX_SEND} and \code{RX_ACK} require \textbf{bit 0 to be set}; the other bits are
    ignored.
  \item \code{ERROR_FLAGS} is cleared \textbf{only by an all-zero word}. \code{0x1} does nothing,
    and \code{0x000000FF} does nothing even though bit 0 is set. The opposite rule to the two
    above.
  \item \code{TX_ID} masks to eleven bits, \code{TX_DLC} to four. A write of \code{0xFFFFFFFF}
    reads back as \code{0x000007FF} and \code{0x0000000F} respectively.
  \item Writes to \code{STATUS} and \code{RX_*} are silently ignored.
  \item Reads of \code{TX_SEND} and \code{RX_ACK} give \code{0x00000000}.
\end{itemize}

\subsection{Data byte order}

\code{TX_DATA} and \code{RX_DATA} are 64 bits wide, with \textbf{data byte 0 in the most
significant bits}:

\begin{console}
bit 63                                                              bit 0
+--------+--------+--------+--------+--------+--------+--------+--------+
| byte 0 | byte 1 | byte 2 | byte 3 | byte 4 | byte 5 | byte 6 | byte 7 |
+--------+--------+--------+--------+--------+--------+--------+--------+
 \___________ TX_DATA_LO ___________/ \___________ TX_DATA_HI __________/
\end{console}

\code{LO} and \code{HI} refer to the order of the register pair, not to bit significance:
\textbf{\code{TX_DATA_LO} carries the first four data bytes.} That is the most common misreading
in the project, and it does not show until a frame arrives back to front.

A frame with \code{dlc = 2} and data \code{AA BB} is therefore packed as
\code{TX_DATA_LO = 0xAABB0000} and \code{TX_DATA_HI = 0x00000000}.

\subsection{The three limitations}
\label{sec:limits}

The hardware is simplified compared with real CAN in many respects. Three of them reach all the
way up into your code, and have to be handled rather than hidden.

\begin{limitation}
\textbf{A lost arbitration looks like a successful transmission.} \code{STATUS} bit 0 comes back
in both cases. The bit means \emph{the transmit port is free again}, not \emph{the frame arrived}.
If you want to know the outcome: wait for bit 0, and \textbf{then} read \code{ERROR_FLAGS}.
\end{limitation}

\begin{limitation}
\textbf{\code{ERROR_FLAGS} is one bit for four causes} --- lost arbitration, stuff error, failed
CRC and a received remote frame. You can see \emph{that} something went wrong, never \emph{what}.
Retry logic that tells them apart cannot be written against this interface.
\end{limitation}

\begin{limitation}
\textbf{No buffering on the receive side.} The controller holds exactly one frame. If the next one
arrives before the previous has been acknowledged it is overwritten, and the newest wins. Poll
often --- and know that at some traffic intensity that will still not be enough.
\end{limitation}

\begin{aside}
\note The hardware validates nothing. A \code{TX_DLC} of \code{0xF} is stored and passed on
exactly as it is. Rejecting invalid frames is \textbf{the driver's} responsibility, and only the
driver's: two copies of the same rule are two things that have to be kept in step.
\end{aside}
